\chapter{РАЗРАБОТКА ПРОГРАММНОГО ОБЕСПЕЧЕНИЯ И АЛГОРИТМОВ РАБОТЫ СИСТЕМЫ}

\section{Структура программного обеспечения микроконтроллера}

Программное обеспечение (ПО) узла мониторинга разработано на языке программирования С++ с использованием среды разработки Arduino IDE и фреймворка ESP-IDF. Архитектура ПО построена по модульному принципу, что обеспечивает удобство отладки и модификации отдельных функций.

Основными программными модулями являются:
\begin{itemize}
    \item модуль инициализации периферии (I2C, GPIO);
    \item модуль управления беспроводным соединением (WiFi Manager);
    \item модуль опроса датчиков и фильтрации данных;
    \item модуль формирования и отправки MQTT-пакетов.
\end{itemize}

\section{Разработка алгоритма функционирования устройства}

Основной цикл работы устройства реализован с использованием механизмов энергосбережения. Алгоритм предусматривает переход микроконтроллера в режим глубокого сна (Deep Sleep) между сеансами передачи данных.



Описание работы алгоритма:
\begin{enumerate}
    \item пробуждение по таймеру;
    \item инициализация датчика SHT31;
    \item чтение данных о температуре и влажности;
    \item проверка условий критического отклонения параметров;
    \item активация модуля Wi-Fi и подключение к брокеру;
    \item публикация данных в топик «sensor/data»;
    \item переход в режим сна на заданный интервал (например, 10 минут).
\end{enumerate}

\section{Разработка структуры пакета данных и протокола обмена}

Для обеспечения совместимости с облачными сервисами визуализации (например, ThingsBoard или Grafana), данные передаются в формате JSON. Пример структуры пакета представлен ниже:

\begin{verbatim}
{
  "dev_id": "ESP32_SRV_01",
  "temp": 24.5,
  "hum": 45.2,
  "rssi": -65,
  "batt": 3.95
}
\end{verbatim}

Использование формата JSON позволяет легко добавлять новые параметры (например, давление или уровень освещенности) без изменения логики работы серверного ПО.

\section{Обработка исключительных ситуаций и ошибок}

Для повышения надежности системы реализован механизм Watchdog Timer (WDT). Если в течение 30 секунд микроконтроллер не завершил цикл передачи данных (например, из-за зависания сетевого стека), происходит автоматический программный сброс устройства.

Также предусмотрено локальное кэширование данных во внутреннюю память Flash при временном отсутствии Wi-Fi соединения. После восстановления связи накопленные данные отправляются на сервер с метками времени.

\newpage